% Part: first-order-logic
% Chapter: model-theory
% Section: theory-of-m

\documentclass[../../../include/open-logic-section]{subfiles}

\begin{document}

% SEGMENT OLP-0188-S01

\olfileid{mod}{bas}{thm}

\section{ஒரு \printtoken{S}{structure} இன் கோட்பாடு}

ஒவ்வோர் !!{structure}~$\Struct{M}$ உம் சில !!{sentence}s ஐ
மெய்யாக்கி, சிலவற்றைப் பொய்யாக்குகிறது. அது மெய்யாக்கும் எல்லா
!!{sentence}s இன் கணம் அதன் \emph{கோட்பாடு} எனப்படுகிறது.
அந்தக் கணம் உண்மையில் ஒரு கோட்பாடே; ஏனெனில் அதிலிருந்து தருக்கமுறையாக
உட்கிடைக்கும் எதுவும்~$\Struct{M}$ உட்பட அதன் எல்லா மாதிரிகளிலும்
மெய்யாக வேண்டும்.

\begin{defn}
  !!a{structure}~$\Struct M$ கொடுக்கப்பட்டால், $\Struct{M}$ இன்
  \emph{கோட்பாடு} என்பது $\Struct{M}$ இல் மெய்யான !!{sentence}s இன்
  கணம் $\Theory{M}$; அதாவது $\Theory{M} =
  \Setabs{!A}{\Sat{M}{!A}}$.
\end{defn}

உய்த்தறிவாக மூடியிருந்தாலும் இல்லாவிட்டாலும், நோக்கமாகக் கொண்ட
பொருள்கோளைக் கொண்ட !!{sentence}s இன் கணங்களைச் சுட்டவும் ``கோட்பாடு''
என்ற சொல்லை முறைசாராமல் பயன்படுத்துகிறோம்.

\begin{prop}
எந்த $\Struct{M}$ க்கும், $\Theory{M}$ நிறைவானது.
\end{prop}

\begin{proof}
எந்த !!{sentence}~$!A$ க்கும், $\Sat{M}{!A}$ அல்லது
$\Sat{M}{\lnot !A}$; எனவே $!A \in \Theory{M}$ அல்லது
$\lnot !A \in \Theory{M}$.
\end{proof}

\begin{prop}\ollabel{prop:equiv}
  ஒவ்வொரு $!A \in \Theory{M}$ க்கும் $\Struct{N} \models !A$
  எனில், $\Struct{M} \elemequiv \Struct{N}$.
\end{prop}

\begin{proof}
எல்லா $!A \in \Theory{M}$ க்கும் $\Sat{N}{!A}$ என்பதால்,
$\Theory{M} \subseteq
\Theory{N}$. $\Sat{N}{!A}$ எனில், $\Sat/{N}{\lnot !A}$; எனவே
$\lnot !A
\notin \Theory{M}$. $\Theory{M}$ நிறைவானது என்பதால், $!A \in
\Theory{M}$. ஆகவே, $\Theory{N} \subseteq \Theory{M}$; எனவே
$\Struct{M} \elemequiv \Struct{N}$.
\end{proof}

\begin{rem}\ollabel{remark:R}
  $\Struct{R} = \langle\Real, <\rangle$ ஐக் கருதுக. இது, மெய்யெண்களின்
  கணம் $\Real$ ஐ ஆட்களமாகக் கொண்ட !!{structure}; மெய்யெண்கள் மீதான
  $<$ உறவாகப் பொருள்கொள்ளப்படும் ஒரே ஒரு 2-இட !!{predicate} ஐ மட்டும்
  கொண்ட !!{language} இற்குரியது. தெளிவாக $\Struct{R}$
  !!{nonenumerable}; எனினும், $\Theory{R}$ வெளிப்படையாக முரணின்மையுடையது
  என்பதால், L\"owenheim--Skolem தேற்றப்படி அதற்கு ஓர் !!a{enumerable}
  மாதிரி உண்டு; அதை $\Struct{S}$ என்க. மேலும் \olref{prop:equiv} படி,
  $\Struct{R}
  \equiv \Struct{S}$. அதோடு, $\Struct{R}$ மற்றும் $\Struct{S}$
  ஓரினமல்ல என்பதால், \olref[iso]{thm:isom} இன் மறுதலை பொதுவாகப்
  பொருந்தாது என்பதை இது காட்டுகிறது.
\end{rem}

\end{document}
